%% =====================================================================
%%  Step 4 -- "Separating hyperplane / Exposed face"
%%  Diagram for the Hahn-Banach proof in ExposedFaceLemmas.lean.
%%
%%  Shared geometry.  The supporting hyperplane is the line
%%
%%      H : y = 0.6 x + 0.3      (equivalently  f(x,y) = y - 0.6x = c = 0.3)
%%
%%  so f <= c is the half-plane BELOW the line and f > c the half-plane ABOVE
%%  it; grad f = (-0.6, 1) points into {f > c}.  The polygon is built so that
%%  the edge P1 P2 lies exactly on H and every other vertex lies strictly
%%  below it: hence P1 P2 is the exposed set Omega \cap H, and the contact
%%  point delta_bound is its midpoint.
%%
%%      P1 = (1.0, 0.9)   on H   (0.6*1.0 + 0.3 = 0.9)
%%      P2 = (4.5, 3.0)   on H   (0.6*4.5 + 0.3 = 3.0)
%%      C  = (5.6, 2.3), D = (4.6, 0.6), E = (1.6, 0.2)   strictly below H
%%      delta_bound = (2.75, 1.95)  midpoint of P1 P2, on H
%%      delta       = (2.90, 1.10)  interior point of Omega, f = -1.64 < c
%%
%%  Build: pdflatex separating_hyperplane.tex
%%         pdftoppm -r 300 -png separating_hyperplane.pdf separating_hyperplane
%% =====================================================================
\documentclass[10pt]{article}
\usepackage{amsmath,amssymb}
\usepackage{tikz}
\usetikzlibrary{arrows.meta,calc,patterns}
\usepackage[active,tightpage]{preview}
\pagestyle{empty}
\setlength{\parindent}{0pt}

%% The polytope Omega: the edge P1 P2 is the exposed set, the rest is below H.
\def\omega{(1.0,0.9) -- (4.5,3.0) -- (5.6,2.3) -- (4.6,0.6) -- (1.6,0.2) -- cycle}
\def\face{(1.0,0.9) -- (4.5,3.0)}              % = Omega \cap H
\def\Hline{(-1.2,-0.42) -- (7.0,4.5)}          % H : y = 0.6x + 0.3
\def\deltaB{2.75,1.95}                        % contact point, on H
\def\deltaI{2.90,1.10}                        % interior root witness

\tikzset{
  omega/.style      = {fill=orange!15, draw=orange!70!black, line width=1.5pt},
  face/.style       = {draw=red!78!black, line width=4.5pt, line cap=round},
  hline/.style      = {draw=red!78!black, line width=1.8pt, -{Latex[length=4mm,width=3mm]}},
  hplain/.style     = {draw=red!70!black, line width=1.6pt},
  gt/.style         = {pattern=north east lines, pattern color=red!40, opacity=0.35},
  lte/.style        = {pattern=north east lines, pattern color=blue!45, opacity=0.35},
  interiorpt/.style = {circle, fill=blue!70!black, draw=white, line width=1.1pt, inner sep=3pt},
  contactpt/.style  = {circle, fill=red!80!black, draw=white, line width=1.3pt, inner sep=3.2pt},
  normal/.style     = {draw=red!78!black, line width=2pt, -{Latex[length=4mm,width=3mm]}},
  sub/.style        = {font=\small\sffamily, fill=white, fill opacity=0.95, text opacity=1,
                        rounded corners=2pt, inner sep=3pt},
  tag/.style        = {font=\large\sffamily\bfseries, fill=white, fill opacity=0.92, text opacity=1,
                        rounded corners=3pt, inner sep=4pt},
}

\begin{document}

%% ---------------------------------------------------------------------
%% 1. The exposed face of a supporting hyperplane.
%% ---------------------------------------------------------------------
\begin{preview}
\begin{tikzpicture}[x=1cm,y=1cm, line join=round]
  % The open half-space {f > c}, shaded.
  \path[gt] (-1.2,-0.42) -- (7.0,4.5) -- (7.0,6.2) -- (-1.2,6.2) -- cycle;

  \draw[omega] \omega;
  \draw[hplain] \Hline;
  \draw[face] \face;

  % grad f = (-0.6, 1) points into {f > c}.
  \draw[normal] (0.30,0.48) -- (-0.52,1.85) node[above, sub] {$f$};

  \node[contactpt]  at (\deltaB) {};
  \node[interiorpt] at (\deltaI) {};
  \node[sub, anchor=south west] at (2.90,1.86) {$x_0=\delta_{\mathrm{bound}}\in F$};
  \node[sub, anchor=north east] at (2.60,1.45) {$\delta\in\operatorname{int}\Omega$};

  \node[tag, text=orange!70!black] at (3.35,1.35) {$\Omega$};
  \node[tag, text=red!70!black]   at (5.20,5.05) {$\{f>c\}$};
  \node[tag, text=blue!65!black]  at (-0.85,-1.15) {$\{f\le c\}$};
  \node[sub, anchor=west]         at (5.65,3.15)   {$H:\ f(x)=c$};
  \node[sub, anchor=south]        at (2.75,0.55)   {$F=\Omega\cap H$};
\end{tikzpicture}
\end{preview}

%% ---------------------------------------------------------------------
%% 2. The same picture, annotated with the content of the theorem.
%% ---------------------------------------------------------------------
\begin{preview}
\begin{tikzpicture}[x=1cm,y=1cm, line join=round]
  \path[gt] (-1.2,-0.42) -- (7.0,4.5) -- (7.0,6.2) -- (-1.2,6.2) -- cycle;

  \draw[omega] \omega;
  \draw[hplain] \Hline;
  \draw[face] \face;

  \node[contactpt]  at (\deltaB) {};
  \node[interiorpt] at (\deltaI) {};

  \node[sub, anchor=west]        at (2.92,1.95)  {$\delta_{\mathrm{bound}}$};
  \node[sub, anchor=south]       at (2.75,0.55)  {$F=\Omega\cap H$};
  \node[tag, text=orange!70!black] at (3.35,1.35) {$\Omega$};
  \node[sub, anchor=north]       at (2.75,-0.75)
    {$\Omega\subseteq\{f\le c\}$, \quad $F=\Omega\cap H\ni\delta_{\mathrm{bound}}$, \quad $s\in R(F)$};
\end{tikzpicture}
\end{preview}

%% ---------------------------------------------------------------------
%% 3. The two half-spaces: Hahn-Banach is strict on the interior, the
%%    closure step makes it non-strict on all of Omega.
%% ---------------------------------------------------------------------
\begin{preview}
\begin{tikzpicture}[x=1cm,y=1cm, line join=round]
  % {f <= c} below and {f > c} above the line.
  \path[lte] (-1.2,-0.42) -- (7.0,4.5) -- (7.0,-2.4) -- (-1.2,-2.4) -- cycle;
  \path[gt]  (-1.2,-0.42) -- (7.0,4.5) -- (7.0,6.2)  -- (-1.2,6.2)  -- cycle;

  \draw[omega] \omega;
  \draw[hplain] \Hline;
  \draw[face] \face;
  \draw[normal] (0.30,0.48) -- (-0.52,1.85) node[above, sub] {$f$};

  \node[contactpt]  at (\deltaB) {};
  \node[interiorpt] at (\deltaI) {};

  \node[tag, text=red!70!black]  at (5.20,5.05)  {$\{f>c\}$};
  \node[tag, text=blue!65!black] at (0.60,-1.85) {$\{f\le c\}$};
  \node[tag, text=orange!70!black] at (3.35,1.35) {$\Omega$};
  \node[sub, anchor=west]       at (2.92,1.95)  {$\delta_{\mathrm{bound}}\in H$};
  \node[sub, anchor=north east] at (2.60,1.45)  {$\delta\notin H$};
  \node[sub, anchor=south]      at (2.75,0.55)  {$F=\Omega\cap H$};
  \node[sub, anchor=north]      at (2.75,-2.95)
    {$f<c$ on $\operatorname{int}\Omega$ \quad lifts to \quad $f\le c$ on $\Omega$, \quad $f(\delta_{\mathrm{bound}})=c$};
\end{tikzpicture}
\end{preview}

\end{document}